\documentclass{article}
\usepackage[utf8]{inputenc}
\usepackage[margin=1.2in]{geometry}
\usepackage{hyperref}

\usepackage{listings}
\usepackage{xcolor}

\definecolor{codegreen}{rgb}{0,0.6,0}
\definecolor{codegray}{rgb}{0.5,0.5,0.5}
\definecolor{codepurple}{rgb}{0.58,0,0.82}
\definecolor{backcolour}{rgb}{0.95,0.95,0.92}

\lstdefinestyle{mystyle}{
    backgroundcolor=\color{backcolour},   
    commentstyle=\color{codegreen},
    keywordstyle=\color{magenta},
    numberstyle=\tiny\color{codegray},
    stringstyle=\color{codepurple},
    basicstyle=\ttfamily\footnotesize,
    breakatwhitespace=false,         
    breaklines=true,                 
    captionpos=b,                    
    keepspaces=true,                 
    numbers=left,                    
    numbersep=5pt,                  
    showspaces=false,                
    showstringspaces=false,
    showtabs=false,                  
    tabsize=2
}

\lstset{style=mystyle}


\usepackage{tikz}
\usetikzlibrary{positioning}

\usepackage{natbib}
\usepackage{graphicx}
\usepackage{amsmath}
\usepackage{amsmath}
\usepackage{amssymb}
\newcommand{\x}{\mathbf{x}}

\title{\vspace{-2 cm}Universidade Federal de Ouro Preto \\ BCC405 \\ Exame Especial}
\author{Prof. Rodrigo Silva}
%\date{}


\begin{document}

\maketitle

% \begin{itemize}
%     \item Implementar as atividades práticas em C++ é altamente recomendado.
%     \item Utilize ao máximo os algoritmos e estruturas de dados da biblioteca STL. \url{https://www.geeksforgeeks.org/the-c-standard-template-library-stl/}. 
%     \item Evite ao máximo a utilização de ponteiros, mas se precisar, utilizar ponteiros inteligentes \url{https://alandefreitas.github.io/moderncpp/basic-syntax/pointers/smart-pointers/}. 
%     \item Quando precisar de uma estrutura de dados linear sempre avalie primeiro a utilização da classe \texttt{vector} (\url{https://en.cppreference.com/w/cpp/container/vector})
% \end{itemize}



\begin{enumerate}
    \item Dado o gradiente da função objetivo, $\nabla f(\mathbf{x})$, defina matematicamente uma direção de descida. 
    \item Dado o gradiente de uma função de restrição de igualdade, $\nabla h(\mathbf{x})$, defina matematicamente uma direção factível.
    \item Dado o gradiente de uma função de restrição de desigualdade, $\nabla g(\mathbf{x})$, defina matematicamente uma direção factível.
    \item Derive a expressão de atualização de $\mathbf{x}$ do método de segunda ordem de Newton utilizando Série de Taylor.
    \item Considere o método de penalidade interior. O que precisa acontecer com $\lambda$ (peso da penalidade) para que o método converja para ótimo restrito? E para o método de penalidade exterior? Justifique em ambos os casos. 
    \item Derive as condições de Karush-Khun-Tucker pela função lagrangeana. 
\end{enumerate}



%\footnotetext{Livro - \textit{Introduction to the Design and Analysis of Algorithms (3rd Edition)}}

%\bibliographystyle{plain}
%\bibliography{references}
\end{document}

